\section{Literature Review}

\subsection{Collaborative filtering and matrix factorization}

Early recommender systems relied on \textbf{collaborative filtering}: inferring preferences from
patterns in the user--item matrix.
Matrix factorization and latent-factor models remain strong baselines when
interaction density is moderate \cite{rendle2009bpr}.

\subsection{Hybrid and content-aware approaches}

\textbf{Hybrid} systems combine collaborative signals with content features (genres, demographics, release year).
Engineered aggregates---such as per-user average rating and per-movie popularity---inject side information
without requiring dense pairwise co-occurrence.
Tree ensembles (random forests, gradient boosting) are effective on such tabular data
\cite{breiman2001random,friedman2001gradient}.

\subsection{Neural collaborative filtering}

He et al.\ \cite{he2017ncf} proposed \textbf{Neural Collaborative Filtering (NCF)}, which learns user and item
embeddings and passes their interaction through a multilayer perceptron to predict ratings.
NCF can capture nonlinear interactions but typically needs more data and tuning than strong gradient-boosted
baselines on medium-scale explicit-feedback datasets.

\subsection{Hyperparameter optimization}

Modern DL pipelines often use automated search.
\textbf{Optuna} provides efficient trial-based optimization for neural architectures and learning rates
\cite{akiba2019optuna}.
We use Optuna for NCF and grid search for sklearn models, selecting configurations on the validation split
before reporting test metrics.

\subsection{Implementation stack}

Modeling uses \textbf{scikit-learn} \cite{pedregosa2011sklearn} for classical ML and
\textbf{PyTorch} \cite{paszke2019pytorch} for NCF.
The MovieLens 1M schema and separator conventions follow the official GroupLens release
\cite{harper2015movielens}.
